\documentclass[twoside]{article}
\usepackage[reviewing]{aag10} % for reviewing
%\usepackage[reviewing,debug]{aag10}\def\extrapage{3} % to produce guidelines
%\usepackage[final]{aag10} % for the final version
\usepackage{graphicx} % for graphics


\begin{document}


\title{Author Guidelines for the Advances in Architectural
	Geometry Conference, London 2014}

\author{Alexander Schiftner \affiliation{Evolute and TU Wien}
	Johannes Wallner \affiliation{TU Graz and TU Wien}}

\maketitle

\def\shortauthor{A. Schiftner and J.\ Wallner}  % for the running head
\def\shorttitle{Author Guidelines for AAG}   % for the running head


\begin{abstract} These author guidelines take care of two things: (i) 
Reviewers should be able to process all submissions in more or less the 
same way, and they should get an impression on how a paper will look in 
print. (ii) The accepted papers must be typeset in the same style, and 
must not contain low quality graphics. 

The submitters are kindly asked to 
take the detailed instructions into account when preparing their 
submission, and in particularly when preparing the camera-ready version of 
their accepted paper. The easiest way to do that is to take the ``Word'' 
and ``\LaTeX'' versions of this document as a template for their own paper. 
\end{abstract}

\section{Introduction}

The {\em Advances in Architectural Geometry} conference welcomes both 
short and long papers, which are being reviewed by the scientific 
committee, possibly by additional external reviewers, under the direction 
of the scientific chairs who make the decision to accept a paper for the 
AAG14 conference volume. In addition the paper can be invited for a 
plenary presentation at one of the two conference days; the other accepted 
papers are invited to present themselves in form of a poster.

The proceedings volume is produced from a camera-ready PDF file provided 
by the authors, and it is imperative that all papers use the same 
typefaces, font sizes, style of headings, size of print area, and so on. 
Therefore we ask you to follow the steps outlined in this document 
carefully when submitting your final manuscript. It is recommended to 
follow the guidelines even for the first submission, because the reviewers 
surely want to know how your paper will look like. You may as well use the 
\LaTeX\ source code or the {\em Microsoft Word} version of this document 
as a template to produce your paper.

Submission for reviewing must be done in in PDF format. The organizers are 
confident that each author has access to software capable of producing PDF 
files (very often you can simply print your document to create a PDF 
file). If you honestly believe that you don't, please contact a 
knowledgeable colleague first, and if this fails, contact the organizers.

\section{Instructions}

The language of both conference and proceedings is English.

\subsection{Margins and page numbering.}

Printed material, including text, illustrations, and charts, but excluding 
the page heading, must be kept within a print area 11.7\,cm (4.6 inches) 
wide by 18.9\,cm (7.44 inches) high. The top margin, i.e., the distance 
between top of the page heading and top of main text, is 12\,mm. You can 
use page numbers for reviewing, but you must not have page numbers for the 
final version.

\subsection{Typefaces and Fonts.}

Wherever {\em Times} is specified, {\em Times Roman} may also be used. If 
neither is available on your word processor, please use the font closest 
in appearance to Times that you have access to (Make sure you use Type-1 
fonts or True-Type fonts).

The title should be in Times 14-point, boldface type, and left-justified. 
Capitalize the first letter of nouns, pronouns, verbs, adjectives, and 
adverbs; do not capitalize articles, coordinate conjunctions, or 
prepositions (unless the title begins with such a word). Leave two blank 
lines after the title.

The Author names and affiliations come beneath the title and are printed 
in Times 10-point (authors in boldface, affiliations non-boldface). Please 
leave two blank lines between title and authors, a half blank line between 
author and affiliation, and a blank line between affiliation and the next 
author.

The abstract should be fully-justified italicized text, below the author 
and affiliation information. Use the word ``Abstract'' as the title, in 
10-point Times, boldface type, left-aligned to the text, initially 
capitalized.

The main text is in 10-point Times, single-spaced. Do \emph{not} use 
double-spacing. All paragraphs should be indented 5\,mm. Make sure your 
text is fully justified -- that is, flush left and flush right. Please do 
not place any additional blank lines between paragraphs. Figure and table 
captions should be 9-point Times type. Page headings, starting with the 
second page, consist of the author names on odd pages, and the title on 
even pages. They are typeset in 9-point Times.  Avoid footnotes if 
possible.\footnote{they should be typeset in 8-point Times.}


\subsection{Section headings.} First and second order headings can be 
numbered in the style seen here. Please do not number third order 
headings.

\paragraph{First order headings.}

They are in 12pt Times, bold-face type. They should be initially 
capitalized, flush left, with one blank line before, and one blank line 
after.

\paragraph{Second order headings.}

They should be 10pt Times, otherwise same as the first order headings.

\paragraph{Third order headings.}

If you really need them, they should be typeset in italics, with one blank 
line before. They are followed by a period and the text.


\begin{figure}[b]
	\includegraphics[width=.5\textwidth]{images/image1}
	\hfill
	\begin{minipage}[b]{.45\textwidth}
\caption{A figure caption can also be located beside the figure
(aligned with the bottom of the figure). It is typeset in 9 pt Times.}
	\end{minipage}
\end{figure}

\begin{figure}[p]
	\centering
	\includegraphics[width=.24\textwidth]{images/image0}
	\includegraphics[width=.24\textwidth]{images/image0}
	\includegraphics[width=.24\textwidth]{images/image0}
	\includegraphics[width=.24\textwidth]{images/image0}
	\includegraphics[width=.30\textwidth]{images/image0}
	\includegraphics[width=.30\textwidth]{images/image0}
	\includegraphics[width=.30\textwidth]{images/image0}
	\includegraphics[width=.40\textwidth]{images/image0}
	\includegraphics[width=.40\textwidth]{images/image0}
	\includegraphics[width=.50\textwidth]{images/image0}
	\caption{Sample figure with short centered caption.}
\end{figure}

\begin{figure}[p]
	\includegraphics[width=\textwidth]{images/image1}
	\caption{Sample figure with a longer caption which then is
typeset in fully justified mode (both left and right). Figure captions 
should be typeset in 9 point Times. Contrary to the instructions found in 
this text, the present figure does not have sufficient resolution to be 
printed this big.}
 \end{figure}

\begin{figure}[p]
	\fboxsep.1mm\fboxrule.1mm
	(a)\hfill\fbox{\includegraphics[width=.4\textwidth]{images/hyp-good}}
	\hfill
	(b)\hfill\fbox{\includegraphics[width=.4\textwidth]{images/hyp-bad}}
	\caption{Mangnified screenshots of PDF files.
{\em Left:} Correctly assembled PDF file. Only minimal artifacts are
visible in areas of high contrast.
{\em Right:} Incorrectly assembled PDF file. Image compression has
not been switched off, artifacts are clearly visible.
The image is not fit for printing.}
	\label{fig:artifacts}
\end{figure}


\subsection{Mathematical symbols.}

The usual guidelines apply: Symbols which designate a single entity (like 
$A$ for the area, $n$ for a number, etc.) should be typeset in italics, 
especially if they appear in the running text. Centered formulas may be 
used for emphasis.
	$$ f(x)=\sin x \Longrightarrow f(90^\circ)=1. $$
 The standard functions (like $\sin$, $\cos$, etc.) are not italicized.


\subsection{References.}

List all bibliographical references at the end of the paper, in 
alphabetical order, with author's names in small capitals and book titles, 
journal names in italics (see the sample bibliography at the end of this 
text). Referencing is done by author(s) and year enclosed in square 
brackets:  \cite{dandelin-1822}, \cite{chusequin:2002}, 
\cite{pauly-2008-sr}.

\subsection{Illustrations, graphs, and photographs}

The proceedings will be printed in color. Usually graphics should be 
centered. If your paper includes images, it is very important that they 
are of sufficient resolution to be faithfully reproduced. To determine the 
optimum size (width and height) of an image, measure the image's size as 
it appears in your document (in millimeters), and then multiply those two 
values by 12. The resulting values are the optimum $x$ and $y$ resolution, 
in pixels, of the image. Image quality will suffer if these guidelines are 
not followed.

{\em Example 1:} An image measures 50\,mm by 75\,mm when placed in a 
document. This image should have a resolution of no less than 600 pixels 
by 900 pixels in order to be reproduced faithfully.

{\em Example 2:} Capturing a screenshot of your entire $1024 \times 768$ 
pixel display monitor may be useful in illustrating a concept from your 
research. In order to be reproduced faithfully, that $1024 \times 768$ 
image should be no larger than 85\,mm by 64\,mm (approximately) when 
placed in your document.

Make sure that the graphics you include in the paper have not been 
compressed or downsampled (as is usually the case with JPEG). Most image 
processing software has a way of specifying the exact amount of 
compression when saving an image, but the default value usually is too 
much. If your graphics has sharp edges, you must not use compression below 
90\%. Otherwise, the compression will be visible in print, such as 
illustrated by Figure \ref{fig:artifacts}b.

\subsection{Hyphenation and typesetting of URLs.}

Some word processors do hyphenation automatically and some don't. The
text becomes prettier if you use hyphens to break long words, because
it keeps whitespace between words smaller. For your final accepted paper
please introduce hyphens manually if necessary. A notorious problem in this
context are URLs (references to the world wide web) which are often treated
as a single word. 

All URLs should be typeset in italics, like
{\http:// www. architecturalgeometry. at/aag14}. \LaTeX\ users
please have a look at the \verb|\http| macro you find in the template file.

% LaTeX users: Note the spaces which occur in the \http command.

\section{Making PDF Files}

When accepted, the final paper must be delivered as a PDF document which 
is not change protected, with all typefaces embedded, and graphics not 
downsampled/compressed. Different ways to achieve this are in more detail 
described below. 

{\it Note:} This part of the guidelines is taken without change from
the guidelines for the 2010 conference. The software described below
may have newer versions available by now.


\subsection{Checking if a PDF file is acceptable.}

In order to check if the final PDF file intended for publication is 
acceptable, please view it in Acrobat Reader, and have a look at the {\em 
List of Fonts used in the document} (see security settings / document 
properties). Each font must be {\em embedded}, and must be a {\em Type-1} 
font or a {\em True-Type} font.

Zoom in on the graphics and check if they exhibit JPEG artifacts (see 
Fig.\ \ref{fig:artifacts}). Artifacts either are already contained in your 
original images or have been introduced by not producing the PDF file 
properly.


\subsection{FreePDF.}


This tool is freely available and known to produce acceptable pdfs. It
creates a new printer device called  ``FreePDF'' which you can print to from
{\it Microsoft Word} or any other word processor. In order to install and use it
please follow this procedure:

$\bullet$  Download the file {\it gs864w32.exe} from
{\http:// sourceforge. net / projects / ghostscript / files/}
(see the folder GPL Ghostscript / 8.64 on the website). Double-click
it and install ghostscript into the default location.

$\bullet$ Download and install {\it FreePDF 4.02} from
{\http://freepdfxp. de/ xpDownload. html}.
 When installing, {\it FreePDF} will find {\it ghostscript} in its
default location.

$\bullet$ Print to the printer {\it FreePDF}. After printing the graphical user
interface of {\it FreePDF} will pop up. Choose the PDF Profile
``High Quality'' and click ``Edit Configuration''. Go to ``Profiles'',
again choose the profile ``High Quality'' and click ``Change''.
Select ``PDF Pictures'' and set all image downsampling options to false.
Click ``Save'', close the configuration dialogue, and save the pdf file.

\subsection{Using word processors with Acrobat Distiller.}

We recommend to download and install the version of the ``CMW'' Adobe 
Acrobat Distiller job options file appropriate for your operating system 
and version of Acrobat from {\http:// www. cadmusmediaworks. com/ 
index2.html} in the ``(Operating System)/\hskip 0pt Applications/\hskip 
0pt Distiller Settings'' folder. The ``CMW'' job options file embeds all 
typefaces and does not downsample or subsample images when creating the 
PDF document.



\subsection{Using \TeX\ and \LaTeX.}

 \LaTeX{} users can use {\em dvips} and {\em ps2pdf} to create a PDF 
document, where {\em Adobe Acrobat Distiller} may be used in place of {\em 
ps2pdf}. In case you don't have access to either please contact the 
organizers once your paper is accepted.

\paragraph{First step: Creating a Postscript file.} {\em dvips} should be 
invoked with the appropriate flags in order to use Type-1 PostScript 
typefaces (if you do not invoke {\em dvips} directly from the command 
line, but by pressing buttons in some editor, make sure to edit your 
configuration and add the flags there).

{\footnotesize\begin{verbatim}
dvips -t a4 -Ppdf -G0 -o myfile.ps myfile.dvi
\end{verbatim}}

\paragraph{Second step: Creating a PDF file.} If you are using version 7.x 
of GhostScript, please use the following method of invoking {\em ps2pdf}, 
in order to embed all typefaces and ensure that images are not downsampled 
or subsampled in the PDF creation process:

{\footnotesize\begin{verbatim}
ps2pdf -dMaxSubsetPct=100 -dCompatibilityLevel=1.3 \
  -dSubsetFonts=true -dEmbedAllFonts=true \
  -dAutoFilterColorImages=false -dAutoFilterGrayImages=false \
  -dColorImageFilter=/FlateEncode -dGrayImageFilter=/FlateEncode \
  -dMonoImageFilter=/FlateEncode mypaper.ps mypaper.pdf
\end{verbatim}}


\noindent If you are using version 8.x of GhostScript, please use this 
method in place of the example above:

{\footnotesize\begin{verbatim}
ps2pdf -dPDFSETTINGS=/prepress -dCompatibilityLevel=1.3 \
  -dAutoFilterColorImages=false -dAutoFilterGrayImages=false \
  -dColorImageFilter=/FlateEncode -dGrayImageFilter=/FlateEncode \
  -dMonoImageFilter=/FlateEncode -dDownsampleColorImages=false \
  -dDownsampleGrayImages=false mypaper.ps mypaper.pdf
\end{verbatim}}


\subsection{Using pdftex/pdflatex.}

{\em pdftex} and {\em pdflatex} (and its variants) can be used only if the 
author can make certain that all typefaces are embedded and images are not 
downsampled or subsampled during the PDF creation process. Configuration 
of these tools to embed all typefaces can be accomplished by editing the 
{\em updmap.cfg} file to enable inclusion of the standard (or base) 14 
typefaces. Linux users can run the {\em updmap} script to do this:

{\footnotesize\begin{verbatim} 
updmap --setoption pdftexDownloadBase14 true 
\end{verbatim}}

Windows users should edit the {\em updmap.cfg} files found in their \TeX\ 
installation directories (one or both of the following may be present):

{\footnotesize\begin{verbatim}
INSTALLDIR\texmf\web2c\updmap.cfg
INSTALLDIR\localtexmf\miktex\config\updmap.cfg
\end{verbatim}}

Ensure the value for {\em pdftexDownloadBase14} is {\em true}, and then 
follow the instructions found at {\http://docs. miktex. org/ manual/} to 
update your MikTeX installation.


\section{Upon acceptance}

After acceptance of your paper, the publisher requires you to sign a form 
to the effect that the usual rights are transferred to him, and that you 
have all necessary rights to include your pictures in the paper. Further 
you will be notified if the chairs have selected your paper for plenary 
presentation.

\section*{Acknowledgments}

\ifreviewing
No Acknowledgments in the reviewing phase
\else
We have borrowed shamelessly from the author instructions of Computer 
graphics conferences and journals, in particular from Eurographics and ACM 
SIGGRAPH.
\fi

\nocite{*} % to make sure all entires in the .bib file actually appear
	   % in the bibliography. Usually this is not needed.

\bibliographystyle{acmsiggraph1} % use ACM SIGGRAPH bibliography style
\bibliography{guidelines} % input your .bib file here

\ifreviewing\else\vfill

\noindent{\it Authors' address:} \\[1ex]
Alexander Schiftner (schiftner@evolute.at):
Evolute GmbH / 
Schwindgasse 4/10 --- A 1040 Wien, Austria. \\
Johannes Wallner (j.wallner@tugraz.at):
Technische Universit\"at Graz / 
Kopernikusgasse 24 --- A 8010 Graz, Austria.\fi
\end{document}
